% =============================================================================
% cap4 §4.2 Datasets públicos (enxuto)
% =============================================================================
\section{Datasets públicos}
\label{sec:metodologia-datasets}

A seleção dos datasets adotados segue três critérios: aderência à tarefa de cada fase, representatividade ecológica do contexto sinantrópico brasileiro (Fases~II e~III) e disponibilidade de protocolo padronizado de avaliação. A Tabela~\ref{tab:datasets-avaliacao} sintetiza as fichas técnicas; todos são distribuídos sob \emph{Community Data License Agreement} permissiva ou licença acadêmica, com agregação dos três primeiros no portal LILA~BC~\cite{lila_bc}.

% Tabela 4.1 — Datasets públicos selecionados
\begin{table}[!htb]
\centering
\caption{Datasets públicos selecionados para a avaliação do pipeline.}
\label{tab:datasets-avaliacao}
\small
\begin{tabular}{@{}p{2.6cm}p{1.8cm}p{2.6cm}p{1.8cm}p{3.6cm}@{}}
\toprule
\textbf{Dataset} & \textbf{Volume} & \textbf{Anotação} & \textbf{Licença} & \textbf{Função no estudo} \\
\midrule
\emph{Caltech CT}~\cite{beery2018terra}        & 243\,k img / 21\,cls  & \emph{bbox} + espécie    & CDLA perm. & Detecção (II); classificação (III) \\
\emph{NACTI}~\cite{tabak2019nacti}             & 3{,}7\,M img / 28\,cls & espécie (bbox parcial)  & CDLA perm. & Par crítico \emph{F.\,catus}/\emph{L.\,rufus} (III) \\
\emph{Snapshot Serengeti}~\cite{norouzzadeh2018automatically} & 7{,}1\,M img / 61\,cls & espécie + contagem & CDLA perm. & Descarte de quadros vazios (II) \\
\emph{PetFace}~\cite{shinoda2024petface}       & 257\,k indiv.         & identidade + raça        & acadêmica  & Re-ID por face (IV) \\
\emph{WildlifeReID-10k}~\cite{adam2024wildlifereid10k} & 10\,k indiv. / 33\,esp. & identidade + \emph{split} & acadêmica & Re-ID por flanco (IV) \\
\bottomrule
\end{tabular}
\end{table}


\subsection*{Fases II e III}
O \textbf{Caltech Camera Traps (CCT)}~\cite{beery2018terra} é a referência primária: 243\,k imagens de 140 câmeras com 21 classes ecologicamente próximas da fauna esperada (gambás, guaxinins, esquilos). A divisão treino/teste oficial separa imagens por local de captura, simulando a aplicação de um classificador genérico em ambiente novo. O \textbf{NACTI}~\cite{tabak2019nacti} fornece o par discriminativo crítico \emph{Felis catus} \emph{vs.}\ \emph{Lynx rufus}, examinando a confusão entre felinos domésticos e silvestres de pequeno porte. O \textbf{Snapshot Serengeti}~\cite{norouzzadeh2018automatically}, com cerca de três quartos de quadros vazios, atua como conjunto de teste severo para o descarte de \emph{ghost frames} na Fase~II.

\subsection*{Fase IV}
O \textbf{\modeloPetFace}~\cite{shinoda2024petface} é a referência primária para re-ID por face: 257\,k indivíduos, subconjunto significativo de gatos domésticos, protocolos oficiais \emph{closed-set} e \emph{open-set}. O \textbf{\modeloWildlife}~\cite{adam2024wildlifereid10k} cobre dez mil indivíduos em $\sim$33 espécies, com divisões \emph{time-aware} e \emph{similarity-aware} que controlam vazamento entre treino e teste, e é a base para o subprocesso de re-ID por flanco e padrão de pelagem (método de partes do corpo de \citeauthor{akbar2025bodypart}, \citeyear{akbar2025bodypart}).

\subsection*{Limites de extrapolação}
Nenhum dataset selecionado foi coletado em ambiente urbano cooperativo com câmeras integradas a abrigos comunitários. As Fases~II e~III tampouco contemplam o conjunto exato de fauna não-alvo do Campus~2 --- em particular, o gato-mourisco (\emph{Herpailurus yagouaroundi}) está ausente, embora classes ecologicamente equivalentes (\emph{Lynx rufus}, \emph{Leopardus pardalis}) estejam presentes. Os resultados configuram \textbf{estimativas de limite superior} do desempenho esperado em condições brasileiras, devendo ser lidos como evidência de viabilidade técnica e não como predição operacional. A validação em condições do Campus~2 fica reservada à implementação futura.
